\Needspace{18\baselineskip}
\Article{19}
Thus we have two methods particularly well suited to practical use. The tables below show their essential arrangements.

\expandafter\def\csname HistoricalNoteText11\endcsname{More precisely, perpendicular to the vertical plane corresponding to the value $u=u^0$, that is, the plane containing the magnetic axis in equilibrium when the second needle is absent. In practice, however, the difference can always safely be neglected, both because it is small and because of the very reasoning we started from in Part III of the preceding article.}
\begin{center}
\textit{First method.}

Both the center and the axis of the second needle lie on a line perpendicular to the magnetic meridian\HistoricalNote{11}.

\begin{tabular}{c|c|c|l|l}
& & &center&north pole\\
Deflection&\multicolumn{2}{c|}{Needle position}&toward&toward\\\hline
$u=u'$&$\psi=90^0$&$U=90^0$&east&east\\
$u=u''$&$\psi=90^0$&$U=270^0$&east&west\\
$u=u'''$&$\psi=270^0$&$U=90^0$&west&east\\
$u=u''''$&$\psi=270^0$&$U=270^0$&west&west
\end{tabular}
\end{center}
\SourceLeaf{112}
\expandafter\def\csname HistoricalNoteText12\endcsname{The needles used in these experiments are about $0^{\mathrm m}3$ long. Had we ventured to include the term $R^{-7}$ in the calculations, we would have diminished the certainty rather than increased it.}
\Needspace{13\baselineskip}
\begin{center}
\textit{Second method.}\par
The center of the second needle lies on the magnetic meridian, and its axis is perpendicular to that meridian.
\begin{tabular}{c|cc|cc}
Deflection&\multicolumn{4}{c}{Needle position}\\
&&&Center toward&North pole toward\\\hline
$u=u'$&$\psi=0$&$U=270^0$&north&west\\
$u=u''$&$\psi=0$&$U=270^0$&north&east\\
$u=u'''$&$\psi=180^0$&$U=90^0$&south&west\\
$u=u''''$&$\psi=180^0$&$U=90^0$&south&east
\end{tabular}
\end{center}

If we then set $\frac14(u'-u''+u'''-u'''')=v$ and
\[
\operatorname{tang}v=LR^{-(n+1)}+L'R^{-(n+3)}+L''R^{-(n+5)}+\text{etc.}
\]
we obtain
\[
\begin{aligned}
\text{first method}\qquad L&=\frac{nmM}{mT+\theta}\\
\text{second method}\qquad L&=\frac{mM}{mT+\theta}
\end{aligned}
\]
